\section{Conclusions}
\label{sec:conclusions}

This paper introduced RTXGNN, a self-explainable graph neural network framework 
for real-time financial fraud detection that treats interpretability as a 
first-class design requirement rather than a post-hoc addition. The central 
finding is that intrinsic explainability and predictive accuracy are not in 
tension: when incorporated during training, sparsity-driven interpretability 
constraints act as domain-informed regularization that benefits both objectives 
simultaneously.

Three limitations bound the current results and define the agenda for future 
work. First, evaluation is concentrated on the Elliptic Bitcoin dataset; 
validation on card-present fraud, ACH transfers, and decentralized finance 
settings is necessary before broad deployment claims can be made. Second, no 
user study with practicing fraud analysts has been conducted, leaving open the 
question of whether SEAL's outputs meaningfully support investigative decisions 
in practice. Third, the framework is vulnerable to abrupt concept drift, as 
demonstrated by the sharp performance collapse at time step 42; continual 
learning mechanisms or online retraining protocols are required for sustained 
production robustness.

Promising directions for extending this work include adaptation to 
continuous-time dynamic graphs, integration of causal explanation frameworks 
to move beyond associative masks, federated variants enabling privacy-preserving 
collaboration across institutions, and adversarial robustness evaluation against 
adaptive fraudsters who may exploit known explanation mechanisms.

\paragraph{Code Availability.} The implementation of RTXGNN is publicly 
available at \url{https://github.com/vinhqdang/explainable_GNN_financial_fraud} 
to facilitate reproducibility and future research.